\section{Task and Benchmark}
\label{sec:task}

\subsection{Bird's-Eye Disaster VLN}
An episode is a tuple
\(
e=(I^{\mathrm{pre}},I^{\mathrm{post}},G,s_0,x,g)
\),
where \(I^{\mathrm{pre}}\) and \(I^{\mathrm{post}}\) are georeferenced pre- and post-event images, \(G\) is the geographic transform, \(s_0\) is the initial UAV state, \(x\) is a natural-language instruction, and \(g\) is the annotated target building and damage label.
At time \(t\), the agent receives a crop determined by latitude, longitude, altitude, and camera geometry.
It chooses a planar waypoint, a sensing action, or stop.
The output contains a final position and a damage judgment.

\paragraph{Metrics.}
Strict success rate \(\SR\) counts an episode only when the agent stops within the target tolerance of the designated building.
Semantic success \(\semSR\) is a diagnostic that accepts a building with the requested class; it must not replace strict \(\SR\).
Navigation error \(\NE\) measures final geodesic distance to the target, and \(\SPL\) discounts successful but inefficient paths.
For damage decisions we report accuracy, ECE, Brier score, and negative log-likelihood.
For active sensing we additionally report trigger coverage, post-trigger decision change, uncertainty reduction, action count, altitude change, and accuracy gain per added sensing action.

\subsection{Existing Benchmark Instance}
The current navigation benchmark contains 40 instructions over xBD-derived scenes.
Instructions vary in target damage class, landmark composition, and difficulty.
The existing E1 comparison uses one run per configuration, so it is a pipeline-scale evaluation rather than a definitive estimate.
The E10 calibration sets contain 21,717 test crops and 69,307 crops in a nominal holdout spanning three disaster types.
Historical model training did not strictly exclude all holdout disaster types; consequently, ``holdout'' here describes the stored split, not a certified unseen-disaster test.

\subsection{Leakage-Safe Evaluation Split}
The strict protocol groups data by disaster event before any crop generation.
No tile, building, temporal pair, or disaster identity may cross training, calibration, validation, and test partitions.
Temperature is fitted on a calibration partition whose events are disjoint from the final test events.
All crops from one building remain in one partition.
The final report includes hashes of manifests and checkpoints, event lists, class counts, and duplicate checks.
Our generated strict manifest contains 12 training events, two validation
events, two test events, and three holdout events, with no event overlap.
The corresponding bitemporal building sets contain 143,230, 73,859, 71,293,
and 69,307 examples, respectively.

\paragraph{Evidence-rich reinspection subset.}
The reinspection study uses 40 approved near-target episodes selected from 60
candidates across ten disaster events.
Every episode has one \texttt{major-damage} or \texttt{destroyed} target.
We inspect contact sheets and reject targets near no-data boundaries, outside
the image, or otherwise failing deterministic geometry checks.
This subset is intentionally all-easy and measures whether the agent uses
visible evidence; it is not used to claim long-range navigation ability.

\paragraph{Evaluation status.}
We distinguish \emph{existing evidence}, which is reported in \secref{sec:experiments}, from the \emph{pre-registered strict evaluation} in Appendix~\ref{app:protocol}.
Unknown strict-split results are not estimated or back-filled.
